\section{Disk equations, compactness, and local gluing}
\label{local:section}

We now fix the symplectic form on $T$ once and for all.  The deformation
parameter will enter only through a relative closed two-form; it will not
change the equation for holomorphic curves.  This section constructs the
local analytic objects needed for the finite integrations in the next
section.  Particular care is needed at an unmarked boundary, at a transverse
boundary jump, and when several boundary nodes are smoothed simultaneously.

\subsection{Boundary labels and the relative weight}

Let $L\subset T$ be the embedded spin Lagrangian of
Proposition~\ref{top:brane}, with its
grading modulo four.  Set $J_*$ to be the flat complex structure determined
by $q+ip$.  We use a fixed sufficiently small uniformly tame neighborhood of
$J_*$ in the space of $\omega$-compatible almost complex structures.
All the structures below take values in this neighborhood.

There are countably many graph labels, arranged in layers indexed by
$\N\cup\{0\}$.  Each layer contains one copy of every scalar graph $G_k$,
$k\geq1$, and of $C$.  A label has the form
\begin{equation}\label{local:graph}
 X=\{p=A_Xq+\nabla f_X(q)\},\qquad
 A_X=h_0k\,1_8\quad\hbox{or}\quad S_-,
\end{equation}
where $f_X$ is a smooth periodic real function.  Layers are ordered
increasingly; within each layer the scalar slopes are ordered decreasingly,
followed by $C$.  Write $X\prec Y$ for this order, and $X^+$ for the copy of
the same slope in the next layer.  The order need not be a well-order; all
analytic constructions first use a finite sublist.  Each graph has the
orientation and spin structure from its base frame, trivial local system,
and the integral grading continued from the horizontal plane through
symmetric matrices.  In particular
\[
 \pi_2(T,X)=0,
\]
by the relative homotopy sequence and injectivity of $\pi_1(X)\to\pi_1(T)$.
Thus a graph supports no nonconstant holomorphic disk.

Choose the potentials so that all labels, including $L$, are pairwise
transverse and have no triple intersections.  In layer zero their Hessians
are as small as needed.  Require $f_{X^+}-f_X$ to be a sufficiently small
perfect Morse function; it may be chosen in the open set containing a small
split sum of cosines.  These choices are simultaneous: enumerate the
countably many required conditions, and use transversality of the
first-jet evaluation of a periodic function.  At each step the conditions
already imposed, as well as the perfect Morse condition, are open, and the
remaining transversality and avoidance conditions are residual in an
arbitrarily small open set.  Equivalently one may use the Baire theorem on
the product of these open sets of successive potential differences.

Our strip convention places the earlier boundary condition at height zero
and the later one at height one, with inputs at positive ends.  For graded
planes with positive rotation angles $\alpha_j\in(0,\pi)$, the degree is
\begin{equation}\label{local:degree}
 |x|=8+\vartheta_2-\vartheta_1-\sum_{j=1}^8\alpha_j/\pi.
\end{equation}
For two graph tangent planes this is the number of positive eigenvalues of
the second Hessian matrix minus the first.  To check this, keep their
difference fixed, move the first matrix to zero, and diagonalize the
difference.  Consequently the maximum of $f_{X^+}-f_X$ gives a unique
degree-zero generator
\[
 w_X\colon X\longrightarrow X^+,
\]
and all generators of $G_i\to G_j$, $i>j$, in layer zero have degree zero.
For arrows incident to $L$, degrees are taken modulo four.

Fix a complex symmetric matrix $M$ commuting with $D$, and put
$B=dq^tM\wedge dp$.  For $L$ and for scalar-slope labels choose forms
$\xi_A$ with $d\xi_A=i_A^*B$.  On $L$ these exist by the pullback-kernel
property proved above.  On a scalar graph we take
\begin{equation}\label{local:primitive}
 \xi_X=-(\nabla f_X)^tM\,dq.
\end{equation}
Indeed symmetry of $M$ makes the scalar linear part of $i_X^*B$ vanish,
and differentiation of \eqref{local:primitive} gives its remaining part.
For a polygon with fixed jump limits set
\begin{equation}\label{local:relative-weight}
 B_{\rm rel}(u)=\int u^*B-\sum_{\text{boundary arcs }a}
                         \int_a u^*\xi_{A(a)},\qquad
 \mathrm{wt}(u)=\qnov^{\int u^*\omega}\exp(tB_{\rm rel}(u)).
\end{equation}
Here $t$ is an ordinary formal variable and $\qnov$ is the Novikov monomial;
the coefficient field is specified in Section~\ref{cat:section}.
Only the unweighted specialization $t=0$ is used on words containing a
$C$-slope label.

\begin{lemma}\label{local:weight}
With the jump limits fixed, $B_{\rm rel}$ is invariant under homotopy
through maps with the prescribed boundary lifts.  It is additive on every
nodal splitting.  Thus the weight in \eqref{local:relative-weight} is
multiplicative on trees, coefficient by coefficient in $t$.
\end{lemma}
\begin{proof}
Apply Stokes to a homotopy cylinder.  Its boundary strips lie on the
prescribed labels, where $B=d\xi_A$, and cancel the variation of the arc
integrals.  The endpoint terms vanish because the jump limits are fixed.
On splitting a disk, integrals add over the components; the inserted
constant boundary paths at a node contribute zero.  This also proves the
assertion for refinements that retain relative classes at every vertex.
\end{proof}

\subsection{The area gap and the later-copy comparison}

\begin{lemma}[Uniform gap for a finite list]\label{local:gap}
For a finite transverse list of the labels and $L$, with no triples, there
is $\delta>0$ such that every nonconstant unforced disk or polygon with
these boundary conditions has area at least $\delta$.  The constant is
independent of the number of marked points and of the domain-dependent
compatible structure in the fixed neighborhood of $J_*$.  In particular
fix $\delta_L>0$ for nonconstant pure disks on $L$ with structure $J_*$.
\end{lemma}
\begin{proof}
Choose finitely many sufficiently small target charts.  A chart containing
a pair intersection meets precisely the two local exact sheets, with that
single intersection; elsewhere a sufficiently small chart meets at most
one sheet.  If the image has sufficiently small diameter, it lies in one
of slightly larger such charts.  In the two-sheet case every jump is at
the same point.  Normalize the two primitives to agree there.  Stokes
then gives zero area, even with repeated jumps; the one-sheet case is the
same.  Positivity implies that the map is constant.

Otherwise the connected image has a fixed positive diameter.  Remove
small disjoint neighborhoods of the finitely many pair intersections.
There is a regular image point outside them through which the image
crosses a ball of a fixed radius, meeting at most one local sheet.  Here
is the monotonicity estimate needed for this assertion.  Let $a(r)$ be
area in a subball and let $\ell(r)$ be the length, with multiplicities, of
its level contours.  Coarea and uniform compatibility give
$a'(r)\geq c\ell(r)$ for almost every $r$.  Each contour arc ending on the
local sheet can be closed there by a path whose length is at most a fixed
multiple of the arc length.  The primitive is exact on that sheet, so
this does not alter the Stokes computation.  A local primitive centered
at a point of each resulting loop gives the quadratic bound
$a(r)\leq C\ell(r)^2$.  Hence
\[
 a'(r)\geq cC^{-1/2}\sqrt{a(r)},\qquad a(r)\geq c_0r^2
\]
up to reducing the fixed radius.  All constants use target charts and
uniform compatibility only.  They do not use derivatives in the domain
of the almost complex structure.  This supplies the asserted gap.
\end{proof}

We impose a further smallness requirement on $X^+$.  It is the geometric
reason its eventual action on the new module will be invertible.

\begin{lemma}[Comparison with vertical fibers]\label{local:comparison}
For each $X$, its next copy $Y=X^+$ may be chosen with a number
$0<\epsilon_X<\delta_L$ and the following properties.  There are symplectic
coordinates in a tube about $X$ in which $X$ is horizontal, $Y$ is
$p=\nabla f(q)$ with $\operatorname{osc}f<\epsilon_X$, and each component
of $L$ in that tube meeting $X$ is a vertical patch at a point $q_j$.
Let $F_j$ be the full vertical fiber there.  Then:
\begin{enumerate}[label=(\roman*)]
\item Unforced $X,Y,L$ triangles of area at most $\epsilon_X$ coincide
literally with $X,Y,F_j$ triangles in the corresponding patch; no such
triangle joins different patches.
\item Every $X,Y,F_j$ triangle, with $X,Y$ corner $r$, has area
$f(r)-f(q_j)\leq\operatorname{osc}f$.  In particular the corner $w_X$
gives area $\Delta_j=f(\max)-f(q_j)<\epsilon_X$.
\item There are no nonconstant $X,L$ or $Y,L$ strips of area at most
$\epsilon_X$, and no nonconstant $X,F_j$ or $Y,F_j$ strips of any area.
\end{enumerate}
The same choices work for any of the allowed almost complex data after
pushing those data forward by the fixed symplectic coordinate change.
If $X\cap L=\varnothing$, choose $Y\cap L=\varnothing$ as well.
\end{lemma}
\begin{proof}
The integral shear $p\mapsto p-A_Xq$ and the gradient translation flatten
$X$.  Near a transverse intersection, the other sheet has equation
$q=q_j+\nabla_pg(p)$, with $dg(0)=0$.  A Hamiltonian shear, cut off in that
chart and generated using $g-g(0)$, makes it vertical while fixing the
horizontal sheet.  There are finitely many disjoint such charts.  By
compactness a sufficiently thin tube about the horizontal torus meets
$L$ only in these vertical patches.  A sufficiently close $Y$ remains an
exact graph throughout this transformation.

The crossing estimate of Lemma~\ref{local:gap}, applied away from the two
nearby graph boundaries, gives a fixed positive area cost for leaving
this tube.  It is uniform as $Y$ approaches $X$.  Choose $\epsilon_X$
below this cost and $\delta_L$, and then choose $Y$ with
$\operatorname{osc}f<\epsilon_X$.  A low-area curve stays in the tube and
its third boundary arc stays in one vertical patch.  Conversely, lift a
triangle with third boundary $F_j$ to the cover and use the primitive
$-p^tdq$.  The horizontal and vertical boundary actions vanish, while
the graph action is $-df$.  Stokes gives $f(r)-f(q_j)$.  Thus every such
triangle also stays in the tube and has the required correspondence.
A strip in a single patch has equal endpoint actions and area zero.
For the full vertical fiber the lifted graph and fiber have a unique
intersection, so this last argument applies without an area bound.
\end{proof}

\subsection{Harmonic domain data}

An $m$-disk has a distinguished boundary output followed
counterclockwise by its ordered inputs.  We use four sectors:
\begin{itemize}
\item pure $m_k$, with all boundary on $L$ and smooth input markings;
\item directed polygons with successive labels
$X_0\prec\cdots\prec X_k$, and all inputs and output transverse jumps;
\item mixed polygons with successive graph labels
$X_0\prec\cdots\prec X_b$, then $L$, with graph inputs followed by an
$X_b,L$ input and any number of smooth $L$ inputs;
\item pure $p_k$ disks with one interior output and $k$ cyclically
ordered smooth boundary inputs.  For nonconstant maps $k=0$ is allowed.
\end{itemize}
We count stable maps only.  The pure unary differential also has its
usual unstable term $d$ on differential forms.  Pure disks, including
all $p$ disks, use $J_*$.  In a sector with jumps, ignore all smooth and
auxiliary markings when choosing the equation.

For an ordered jump-arc list $\mathbf A=(A_0,\ldots,A_{m-1})$, let
$h_j(z)$ denote harmonic measure of its $j$th arc as seen from $z$.
Equivalently, send $z$ to the center of the disk and divide the arc's
angular length by $2\pi$.  Put
\begin{equation}\label{local:harmonic-J}
 J(z,x)=J_{\mathbf A}(h_0(z),\ldots,h_{m-1}(z);x),
 \qquad h(z)\in\Delta^{m-1}.
\end{equation}
The functions $J_{\mathbf A}$ are smooth on the closed simplices.

\begin{proposition}\label{local:data}
The functions in \eqref{local:harmonic-J} can be chosen simultaneously
with the following properties.
\begin{enumerate}[label=(\roman*)]
\item On a simplex face they agree with the function for the surviving
subsequence.  They equal $J_*$ near each vertex.
\item Near an edge joining labels $A,B$, they equal exactly
$J_{AB,h_B}$, for a height-dependent strip datum constant $J_*$ near
heights zero and one.  Take $J_{XL,\tau}=J_*$.
\item All nonconstant graph-pair strips and all directed polygons are
regular, including domain variations.  The comparison triangles of
Lemma~\ref{local:comparison}, with $L$ replaced by any $F_j$, are regular
for the pushforward of the same triangle data.
\end{enumerate}
No regularity of general mixed or pure maps is required.
\end{proposition}
\begin{proof}
First prescribe generic graph-pair data, supported in heights
$(0.1,0.9)$.  Ordinary strip transversality applies to these
$s$-independent data.  Its injective-point condition can be expressed as
injectivity of $(s,\tau)\mapsto(\tau,u(s,\tau))$: if an open immersed
patch had a same-height partner on every point, a Baire argument among
compact immersed patches would give overlapping branches
$u(s,\tau)=u(\phi(s,\tau),\tau)$.  The equation forces
$\phi(s,\tau)=s+c$; unique continuation would then contradict finite
energy for a nonconstant strip.  Critical images are negligible and
immersive points are dense.  Hence variations of $J_\tau$ near an
injective point kill any proposed cokernel element, also when the
variation is restricted to the indicated open height interval.

Extend by word size.  Require the pair formula where
$h_A+h_B>c_m$ and the vertex formula where $h_A>d_m$, using thresholds
increasing slightly with $m$, with $c_m$ close to one and $d_m$ close to
$0.94$.  Pair variations vanish near the endpoints, so these requirements
agree on overlaps.  In a local linear chart of compatible structures
near $J_*$, first extend the required collar formulas and then extend
the remaining boundary values over the simplex.  Inclusion--exclusion
in corner charts extends functions agreeing on intersecting faces;
a partition of unity combines the extensions.  A smooth retraction
into the chosen small neighborhood retains all already prescribed
values.

For $m\geq3$, a jump configuration together with one interior point,
modulo disk automorphisms, is identified with the open simplex by $h$:
after sending the point to the center, the successive arc lengths
recover all endpoints up to rotation.  There is a nonempty free open
set outside the prescribed edge and vertex collars.  Each fixed domain
has interior points mapping to this set under $h$: send any three of
its jumps to equally spaced circle points.  At the center every
individual weight is at most $1/3$ and every pair has total at most
$2/3$, so a neighborhood of the center is free.  There is an
immersive point there for a nonconstant solution.  A variation of
$J_{\mathbf A}$ supported near that simplex point isolates its domain
point, even if the map has multiple images.  The usual pairing with a
nonzero adjoint solution and unique continuation prove surjectivity
of the universal linearization.  A directed map is nonconstant because
there are no triples.  Sard--Smale therefore gives the required
regularity relative to all prescribed faces and collars.  Apply the
same argument also to the finitely or countably many auxiliary fiber
triangles using precisely the pushed-forward datum.  There are only
countably many words, classes, and such transversality requirements,
so the successive residual choices can all be made.
\end{proof}

\subsection{Boundary compactification and finite cuts}

\begin{lemma}[Only disk components occur]\label{local:disk-trees}
Every stable limit in the four sectors consists of disks joined at
boundary nodes.  Smooth nodes in a mixed sector split off pure $L$
disks.  At jump nodes the two ends have opposite transitions, and each
vertex has the corresponding subsequence of the original ordered
labels.  The same conclusion holds for positive-area $p_0$ disks;
there is no additional stratum in which the whole boundary collapses.
\end{lemma}
\begin{proof}
There are no nonconstant spheres because $\pi_2(T)=0$.  The disk vertices
in the doubled domain form a connected subtree: they are fixed vertices
of its real involution, and the path between fixed vertices is fixed.
Consequently other sphere trees could only hang from a disk at an
interior node.  A tree of $v$ constant spheres with one attachment and
at most one interior marking has at most $2v$ special points counted
with multiplicity, whereas stability requires at least $3v$; hence no
such tree occurs.

The issue when $p_0$ has no boundary markings is the possible collapse
of the entire real circle.  In the doubled dual tree the collapsed
boundary would be the midpoint of an edge whose branches are exchanged
by the involution.  This cannot coexist with a fixed disk vertex:
the path from a fixed vertex to that midpoint would be fixed, including
its final half-edge, contrary to exchanging the two half-edges.  Thus
all normalized components occur in conjugate sphere pairs.  In the
original half these are sphere components.  Positive area would force
one of them to be nonconstant, which is impossible.  This rules out
boundary collapse in the only sector where it is not excluded by a
boundary marking.

Following labels around a boundary smoothing proves the subsequence
assertion.  A smooth node on a graph would cut off a single-graph
subtree with no original inputs.  An outermost nonconstant component
would be a disk on that weakly exact graph, while a wholly constant
subtree is unstable.  Such a subtree is impossible.  Thus smooth
attachments in a mixed sector are pure $L$ branches.  On a limiting
core the omitted arcs have harmonic measure zero, so
Proposition~\ref{local:data}(i) gives exactly the equation of that vertex.
\end{proof}

We will verify compactness, including absence of neck drift, after giving
the explicit collar coefficients below.  Its finite-cut consequence can
already be stated.  For a finite list and a positive integer $N$ choose
\[
 0<a<\min(1,\delta_L,\delta),
\]
where $\delta$ is the gap of Lemma~\ref{local:gap}.  Define
\begin{equation}\label{local:cost}
 c_m(\beta,k)=\omega(\beta)+a(k-1),\qquad
 c_p(\beta,k)=\omega(\beta)+ak.
\end{equation}
Retain all stable types of cost at most $4N$, with their relative-class
refinements and all nodal types whether or not they are smoothable.

\begin{proposition}\label{local:finite}
This is a finite collection of compact moduli problems, closed under
normalized boundary and corners.  Costs add on every splitting, every
stable vertex has positive cost, and only finitely many relative classes
occur.  No discreteness of the symplectic period group is needed.
\end{proposition}
\begin{proof}
For an $m$ splitting with child arity $s$ and parent arity $k-s+1$,
$(s-1)+(k-s)=k-1$.  For a $p$ root the analogous equality is
$(s-1)+(k-s+1)=k$.  This proves additivity, including $s=0$.
A stable constant pure $m$ vertex has $k\geq2$, and a constant $p$
vertex has $k\geq1$.  A stable constant mixed vertex likewise has at
least three boundary special points; directed constants are excluded
by no triples.  Such vertices cost at least $a$.  Every nonconstant
$m$ vertex costs at least $\delta-a>0$, and every nonconstant $p$
vertex costs at least $\delta$.  Thus the cut bounds the number of
vertices, total area, and arity.  Only finitely many combinatorial
patterns and asymptotic generators occur.

For completeness the compactness argument, proved below in
Lemma~\ref{local:compactness}, is for the union over relative classes
with these bounds.  In a convergent sequence, maps on retained compact
cores are homotopic to their limits, and the remaining thin parts
have arbitrarily small diameter in the appropriate local one- or
two-sheet charts.  Contracting these parts and matching at the nodes
preserves the total relative class.  The class is therefore locally
constant in the compact stable-map space.  A finite subcover yields
finitely many classes.  The same reasoning retains any relative
refinement used to record the additive weight.  Positivity and
additivity show that every component of a boundary type stays in the
cut, proving boundary closure.
\end{proof}

\subsection{Real plumbing and quantitative collar coefficients}

Since all nodes are boundary nodes, we can use whole-disk projective
coordinates for every retained core.  Root a disk tree at its boundary
output, or at the disk carrying the interior output in the $p$ sector.
On a child use an upper-half-plane coordinate $z$ with its output at
infinity; on its parent put the attaching point at $a_e$.  The insertion
\begin{equation}\label{local:plumbing}
 \zeta=a_e+\rho_e z,\qquad \rho_e=e^{-\pi T_e},\qquad T_e\gg1,
\end{equation}
is real plumbing, using $-1/z$ as the child's nodal coordinate.  Cutting
out the corresponding half-disks and making this identification gives
another upper half-plane.  Iteration therefore identifies each core by
the restriction of a whole-disk biholomorphism.  This point will also
make the intrinsic normalizations in the next section compatible with
gluing.

Use node-specific projective coordinates varying smoothly in stabilized
component configurations.  They can be chosen equivariantly: a finite-order
orientation-preserving disk automorphism fixing a boundary point is the
identity, so a node stabilizer acts trivially on its incident disks.
Choose the germs on representatives of the node orbits and transport
them; invariant local slices for the finite stabilizer handle component
configurations.  Descendant arc endpoints in a parent coordinate have
the form $a_e+O(\rho_e)$, and exterior endpoints in a child coordinate
have reciprocals $O(\rho_e)$.  Iterated insertions give smooth expressions
in all plumbing parameters; apparent denominators stay separated by
using the nearest common ancestor coordinate.

\begin{lemma}[Harmonic collar estimates]\label{local:neck-coefficients}
On a jump neck $e$, write
\[
 \zeta-a_e=-e^{-\pi(s+i\tau)},\quad
 R\leq s\leq T_e-R,\quad 0\leq\tau\leq1,
 \qquad r_l=e^{-\pi s},\quad r_r=e^{-\pi(T_e-s)}.
\]
The transitions at the ends are an ordered pair $A,B$.  For sufficiently
large fixed $R$, the equation uses precisely its pair collar, and
\begin{equation}\label{local:angle-estimate}
 h_B=\tau+O(r_l+r_r),\qquad
 h_A=1-\tau+O(r_l+r_r),\qquad
 \sum_{C\ne A,B}h_C=O(r_l+r_r).
\end{equation}
On the parent half the difference from the partially broken parent
coefficient is $O(r_r)$; on the child half the corresponding difference
is $O(r_l)$.  These estimates hold after any prescribed finite number of
strip, target, and compact component-parameter derivatives.  Logarithmic
length derivatives have the same estimates, and a derivative with
respect to another length $T_f$ carries an additional factor
$O(e^{-\pi T_f})$, uniformly in the other large lengths.  In particular
all coefficients have uniformly bounded strip derivatives.

On a smooth matching neck the datum is exactly $J_*$ throughout a long
middle collar.  The same one-sided estimates hold at external jump ends.
\end{lemma}
\begin{proof}
Harmonic measure in the upper half-plane is the angle subtended by an
interval, divided by $\pi$.  For the two intervals adjacent to $a_e$,
the limiting angles in the displayed coordinate are $\pi\tau$ and
$\pi(1-\tau)$.  A displacement of a clustered endpoint is $O(\rho_e)$;
dividing by $|\zeta-a_e|=r_l$ gives $O(r_r)$.  A separated exterior
endpoint contributes an angle $O(r_l)$, with its reciprocal distance
bounded.  The smooth angle function and all of its derivatives on these
separated configurations give \eqref{local:angle-estimate} and the two
one-sided comparisons.  Differentiating $\rho_e$ logarithmically only
multiplies it by $-\pi$.  If an endpoint shift comes from a descendant
node $f$, its differentiated displacement is $O(\rho_e\rho_f)$ before
dividing by $r_l$; the same argument in an ancestor or nearest-common-ancestor chart treats all other cases.  This proves uniformity in every
other length as well as all finite derivative versions.

For large $R$ the total omitted weight is uniformly small, so the
simplex point lies in the pair neighborhood of
Proposition~\ref{local:data}; the coefficient is exactly $J_{AB,h_B}$.
At a smooth node only a pure $L$ branch is inserted.  No original jump
lies down that branch.  On the long collar the weight of $L$ is in the
vertex region, where the datum is exactly $J_*$.
\end{proof}

\subsection{Decay and compactness}

\begin{lemma}[Transverse ends and smooth-node decay]\label{local:decay}
At a jump limit $x\in A\cap B$, the operator
\[
 \mathcal A\xi=J_{AB,\tau}(x)\partial_\tau\xi,\qquad
 \xi(0)\in T_xA,\quad\xi(1)\in T_xB,
\]
on the interval is self-adjoint for the $\tau$-dependent compatible
$L^2$ metric, has compact resolvent, and has zero kernel.  In particular
it has a positive spectral gap.

On a sufficiently small-image long neck, solutions with no obstruction
support satisfy uniform exponential estimates from its two ends for
all positive derivatives.  At a jump the limiting value is the fixed
intersection; at a smooth node one separates a single movable value in
$L$.  On a strip $[0,T]\times[0,1]$, after removing fixed end collars,
one may take
\begin{equation}\label{local:two-sided-decay}
 |\nabla^jdu(s,\tau)|\leq C_j\bigl(e^{-\mu s}+e^{-\mu(T-s)}\bigr)
 \quad (j\leq j_0),
\end{equation}
for any preassigned finite $j_0$, with constants uniform near the given
nodal map.  The same estimates hold for small compact-core-forced
solutions, since their neck equations are unforced.
\end{lemma}
\begin{proof}
Put $g_\tau(v,w)=\omega(v,J_{AB,\tau}(x)w)$.  Compatibility gives
\[
 \int_0^1g_\tau(\mathcal A\xi,\eta)\,d\tau
   =\int_0^1\omega(\partial_\tau\xi,\eta)\,d\tau
   =-\int_0^1\omega(\xi,\partial_\tau\eta)\,d\tau
   =\int_0^1g_\tau(\xi,\mathcal A\eta)\,d\tau.
\]
The boundary term vanishes by Lagrangianity.  The adjoint has the same
boundary domain, which proves self-adjointness; the one-dimensional
elliptic estimate and compact embedding give compact resolvent.
Its kernel consists of constant paths in $T_xA\cap T_xB$, hence is zero.
Notice that no constancy of $J_{AB,\tau}$ in $\tau$ was used.

The full-cylinder inverse of $\partial_s+\mathcal A$ is obtained by
integrating the positive spectral subspace forward and the negative
subspace backward.  Its kernels are bounded by $Ce^{-\mu_0|s-s'|}$,
where $\mu_0$ is the spectral gap.  This gives the weighted estimate for
any logarithmic weight slope less than $\mu_0$; local elliptic estimates
give its Sobolev versions.  Straighten the two boundary sheets at $x$.
The nonlinear equation on a small-image neck differs from this model
by a small first-order term, a zero-order term bounded by $C|du|$, and
the controlled coefficients in Lemma~\ref{local:neck-coefficients}.
After a fixed shortening of the collar, absorb the small first-order
term in the cylinder estimate.  Cutting off at the two ends yields
\eqref{local:two-sided-decay}; differentiation and elliptic bootstrapping
give the asserted finite range of derivatives.

For a smooth node both boundaries lie on the same sheet.  Work first
in its cotangent tube, whose primitive vanishes on the zero section.
For a short path with endpoints on the sheet, its action is bounded
by a fixed constant times the square of its length: close it by a
short sheet path and use a centered primitive in a local chart.  If
$E(R)$ is the energy in $[R,T-R]\times[0,1]$, Stokes and this estimate
at the two slices give
\begin{equation}\label{local:neck-energy}
 E(R)\leq C\bigl(e(R)+e(T-R)\bigr)=-C E'(R),
\end{equation}
where $e(s)$ is slice energy, up to a fixed uniform comparison of
metrics.  Therefore $E(R)\leq C'e^{-R/C}$.  Small-energy regularity on
unit strips, with the uniform coefficient bounds above, gives
exponential decay of positive derivatives.  Integrating the first
derivative gives convergence to one value and the same two-sided
estimate measured relative to a middle boundary value.  This argument
also shows that the constant asymptotic mode must be retained, rather
than estimated by an invertible transverse model.  The argument for
two transverse sheets can alternatively use \eqref{local:neck-energy}
by closing each slice to their intersection along the sheets.
\end{proof}

\begin{lemma}\label{local:compactness}
For a finite list, fixed arity bound, and bounded area, the unforced
stable-map spaces of the four sectors are compact after adjoining the
disk trees of Lemma~\ref{local:disk-trees}.  There is no lost energy or
drift of the matching values.  On the cores the convergence is smooth,
and relative classes agree eventually with the total limiting class.
\end{lemma}
\begin{proof}
Apply the usual disk and strip rescaling procedure, retaining marked-point clustering scales as well as nonconstant bubbling scales.  The
required coefficient compactness holds here directly: harmonic
measures are bounded harmonic functions and have uniform derivatives
on interior hyperbolic balls.  Under any disk reparametrization,
clustered endpoints converge to a limiting configuration; on a
retained core their harmonic measures converge to the surviving ones.
Along a retained boundary arc the datum is constant in a vertex
neighborhood; on a jump collar use
Lemma~\ref{local:neck-coefficients}.  An interior blowup at scale much
smaller than the boundary distance sees constant-in-domain data.
Consequently any new gradient concentration gives an ordinary disk,
a transverse strip, or an interior sphere.  Spheres are impossible,
and each nonconstant level consumes the positive area gap.  There
are therefore only finitely many nonconstant levels, with precisely
the equations prescribed above.

After all these levels have been extracted, derivatives on the
remaining middle strips are small by small-energy regularity.  At
a jump, a short slice between two distinct compact sheets is close
to their finite intersection set.  Continuity makes it remain near
one intersection throughout each remaining thin strip.  For one
sheet, short slices already lie in its cotangent tube, even before
one knows that their longitudinal motion is small.  Thus
\eqref{local:neck-energy} applies in both cases.  The constants in the
single-sheet action estimate are uniform along the compact sheet;
it does not assume the image lies near a fixed point in advance.
Its exponential consequence and integration of $|du|$ show that the
central oscillation tends to zero.  This excludes slow travel along
$L$, as well as loss of energy, and identifies the two limiting values
at a smooth node.  Elliptic estimates give the claimed core
convergence.  Finally, thick-region homotopies and contractions of
these small thin parts within their prescribed local sheets identify
the total relative classes.  Stability and
Lemma~\ref{local:disk-trees} give the stated compactification.
\end{proof}

\subsection{Fixed-domain augmentation}

We specify precisely the local input from the established disk theory.
For a compact target, a compact relatively spin boundary Lagrangian,
and fixed tame $J$, compact-core obstruction data give finite-dimensional
disk charts and smooth evaluations; augmented boundary-evaluation
surjectivity gives an output submersion
\cite[Theorem~7.1, Condition~7.11, and Lemma~7.12]{fooo-i}.
The gluing theorem assumes surjectivity of the augmented component
operators and of the difference of their matching evaluations.
It gives local gluing and exponentially small positive length
derivatives on retained cores, also with domain parameters
\cite[Assumption~3.12, Theorems~3.13, 6.4, and 8.16--8.17]{fooo-decay}.
For a fixed datum our finite union of labels is an immersed Lagrangian
with transverse double points; retaining its boundary-lift sectors
provides the standard transverse-end Fredholm and determinant-line
setup \cite[Sections~4--5]{akaho-joyce}.  We next prove the extensions
needed for \eqref{local:harmonic-J} and for $p_0$.

At a component $u$ use high Sobolev order on its compact part.  At a
jump take weighted sections tending to zero at the fixed intersection.
At a smooth special point viewed as a cylindrical end, write a section
as its movable limit $a\in T_{u(z)}L$ plus a decaying section.  Small
positive weights are chosen below all nonzero indicial roots.  This
is the usual Fredholm formulation of a removable smooth boundary
point: the zero mode is the tangent limit, and homogeneous or
compact-core-forced solutions extend over the point.  Stabilize any
unstable nonconstant domain by interior marks constrained to transverse
target slices at immersive points, taking finite automorphism orbits.
One can do this also for an unbroken $p_0$ map.  The original interior
mark is retained, but no boundary mark is artificially made an output.

\begin{lemma}\label{local:augmentation}
For each component one can choose a finite-dimensional space $E$ and
an insertion $I:E\to C_c^\infty(\Lambda^{0,1}\otimes u^*TT)$ supported
in interior cores away from all markings, such that $D_u\bar\partial-I$
is surjective at fixed domain.  For every pure $m$ component the space
can be chosen so that
\begin{equation}\label{local:augmented-output}
 (\xi,e)\longmapsto
 (D_u\bar\partial\xi-Ie,\,\xi(z_0))
\end{equation}
is surjective as well.  Componentwise sums give a surjective nodal
operator with all smooth matching equations imposed.  The same holds
with a prescribed pure root-output variation.
\end{lemma}
\begin{proof}
The Fredholm cokernel is finite-dimensional.  At first Sobolev order
with target weighted $L^p$, $p>2$, smooth interior tests away from the
finite special points are dense, and hence span its quotient.  Choose
finitely many such tests.  To arrange \eqref{local:augmented-output},
apply the same argument to the operator restricted to sections with
$\xi(z_0)=0$.  Given any prescribed output value, first choose a
section attaining it and then solve its remaining equation with zero
output variation.  Bootstrapping transfers these conclusions to the
higher Sobolev orders.  Replace $E$ by the formal sum of its translates
under the finite stabilizer and average the auxiliary metrics.  This
makes the choice equivariant without reducing its image; injectivity
of the formal insertion is not needed.

At a smooth node the outgoing child is a pure $m$ component.  Starting
at the root, choose its solution with any required external output;
then solve each child with output value equal to the already chosen
parent value, and continue outward.  Formula
\eqref{local:augmented-output} permits this at each step.  Jump
matching is at a fixed transverse intersection and adds no tangent
matching variable.  This proves nodal surjectivity.  The pure constant
disk operator is already surjective with arbitrary constant boundary
value.  A constant mixed component has only an $X,L$ pair at the same
intersection, together with removable smooth marks, and has the
invertible transverse strip operator.  Thus constants require no
additional interior tests.
\end{proof}

\subsection{A uniform gluing estimate with several necks}

We give the uniform estimate explicitly, rather than invoking a fixed-$J$
gluing theorem for the domain-dependent equation.  Fix a stabilized nodal
map and the component augmentations of Lemma~\ref{local:augmentation}.
All constants below are uniform on a small relatively compact neighborhood
of its augmented nodal solutions.  The finitely many component operators
have a common indicial gap.  Choose a small positive weight $\delta$ below
one tenth of that gap and of the decay rates in
Lemmas~\ref{local:neck-coefficients} and \ref{local:decay}.

On an end use weight $e^{\delta s}$.  On a joined neck of length $T_e$
use a smooth version of
\begin{equation}\label{local:weight-function}
 w_e(s)=e^{\delta\min(s,T_e-s)}.
\end{equation}
The smoothing takes place in a fixed-width middle interval and changes
all norms by uniform factors.  At a jump, sections have zero asymptotic
constant.  At a smooth node write
\[
 \xi_e=a_e+\xi_e^{\rm dec},\qquad a_e\in T_{x_e}L,
\]
using fixed sheet-straightening coordinates; include $|a_e|$ separately
in the norm and weight only the remainder.  On a joined neck $a_e$ may
be the value at its middle boundary point.  On the finite set of cores
use the ordinary $W^{m,2}$ norms, and on remainders use
$\sum_{j\leq m}\|w_e\nabla^j\xi_e^{\rm dec}\|_{L^2}$.
Add the finite-dimensional domain and coefficient norms.  Denote the
resulting norm by $\|\cdot\|_{m,\delta,T}$, and use the corresponding
weighted $W^{m-1,2}$ target norm.  Here $m$ can be taken arbitrarily
large, but finite.

\begin{lemma}[Uniform augmented inverse]\label{local:inverse}
Fix any subset of nodes to remain broken.  For all other lengths
$T_e$ sufficiently large, the linearized thickened equation at the
preglued map is surjective, uniformly in every length.  After adding
finite-dimensional kernel coordinates measured on the cores and in
the coefficient variables, it has an inverse $Q_T$ satisfying
\begin{equation}\label{local:inverse-estimate}
 \|Q_T(\eta,b)\|_{m,\delta,T}
       \leq C_m\bigl(\|\eta\|_{m-1,\delta,T}+|b|\bigr).
\end{equation}
The same holds at sufficiently small thickened solutions, at fixed
domain, and with any prescribed pure root-output variation.
\end{lemma}
\begin{proof}
Choose core observables giving an isomorphism from the augmented nodal
kernel to a Euclidean space.  Together with the nodal operator and its
matching equations, these observables form a bounded isomorphism.
The same isomorphism works in the required finite range of Sobolev
orders and small weights: the indicial gaps do not change the index,
and elliptic regularity identifies its kernels.  Domain coordinates
can either be included or held fixed.  The output assertion uses the
augmented operator \eqref{local:augmented-output}.

Here is a bounded parametrix for the smoothed operator.  Split a target
section at each middle neck by a partition whose transition has fixed
width.  Transfer the two pieces to the corresponding semi-infinite
ends and extend them by zero.  Their one-sided weighted norms are
bounded by a fixed multiple of the joined norm, since on the support
of the left piece $s\leq T_e/2+O(1)$ and hence
\[
 e^{\delta s}/w_e(s)\leq e^{O(\delta)};
\]
the right side has the identical estimate.  Solve this forcing with
the nodal isomorphism, retaining all smooth matching equations and
the prescribed core coordinates.  The resulting two sections at a
smooth node have the same constant $a_e$.

On a component end cut off only the decaying part.  Take the left
cutoff equal to one through $T_e/2+2$, zero by $3T_e/4$, and with
$j$th derivative $O(T_e^{-j})$; take the reflected right cutoff.
The support of each split forcing lies where its own cutoff is one.
On the overlap set
\begin{equation}\label{local:splice}
 \xi_e=a_e+\beta_l\xi_l^{\rm dec}
                  +\beta_r\xi_r^{\rm dec}
\end{equation}
at a smooth node; at a jump omit $a_e$.  Thus the matching constant
is included once, not twice.  The joined weight is bounded by either
one-sided weight where that end is used, so the decaying terms have
uniform norm.  If the joined norm uses the actual middle value
instead of $a_e$, weighted Sobolev embedding gives
\[
 |\xi_e(T_e/2,0)-a_e|
       \leq C e^{-\delta T_e/2}
                (\|\xi_l^{\rm dec}\|_{m,\delta}
                 +\|\xi_r^{\rm dec}\|_{m,\delta}).
\]
Since $\|w_e\|_{L^2}\leq Ce^{\delta T_e/2}$, recentering by this
constant is still uniformly bounded.  This verifies the parametrix
bound including the zero modes.

Apply the smoothed operator to \eqref{local:splice}.  Commutators with
the cutoffs have norm $O(T_e^{-1})$.  The difference of the equation
coefficients from the corresponding component coefficients on a
cut end is exponentially small: its cutoff ends at $3T_e/4$, so the
opposite-end errors of Lemma~\ref{local:neck-coefficients} are at most
$Ce^{-\pi T_e/4}$ there, and pregluing errors have the analogous
map-decay bound.  At a smooth match, the limiting operator annihilates
a constant tangent vector.  Its zero-order coefficient on such a
vector is $O(|du|)$, so all errors involving $a_e$ have the same
exponential estimate after weighting; the choice of $\delta$ makes
this estimate summable.  Core-supported obstruction maps, constraints,
and kernel observables converge by their smooth core transport.
It follows that the parametrix $P_T$ satisfies
\begin{equation}\label{local:parametrix-error}
 \|\mathcal L_TP_T-\id\|
   \leq C_m\sum_{e\ \mathrm{smoothed}}
                      (T_e^{-1}+e^{-cT_e}),
\end{equation}
where $\mathcal L_T$ includes the kernel observables.  The number of
nodes is fixed.  Neumann inversion proves a uniformly bounded right
inverse.  The boundary-node Fredholm index formula subtracts
$\dim L$ for a movable smooth match and adds no matching dimension
at a nondegenerate transverse end.  It is exactly the nodal index
with the chosen smoothing lengths held fixed.  Hence adding the
same kernel observables gives index zero, so this right inverse is
an inverse.  All these estimates hold with any other nodes kept
broken.  Small changes of the map and coefficients preserve the
bound by a final Neumann inversion.
\end{proof}

\begin{figure}[ht]
\centering
\begin{tikzpicture}[>=Latex,scale=.91]
 \draw[thick] (0,0)--(11,0);
 \draw (0,-.15)--(0,.15) (5.5,-.15)--(5.5,.15) (11,-.15)--(11,.15);
 \node[above] at (0,.15) {parent core};
 \node[above] at (11,.15) {child core};
 \node[above] at (5.5,.15) {$a_e$ included once};
 \node[below] at (0,-.15) {$s=0$};
 \node[below] at (5.5,-.15) {$s=T_e/2$};
 \node[below] at (11,-.15) {$s=T_e$};
 \draw[->,blue!60!black] (.7,1.2)--(7.8,1.2);
 \node[above,blue!60!black] at (3.5,1.2) {$\beta_l\xi_l^{\rm dec}$};
 \draw[<-,red!60!black] (3.2,-1.15)--(10.3,-1.15);
 \node[below,red!60!black] at (7.5,-1.15) {$\beta_r\xi_r^{\rm dec}$};
\end{tikzpicture}
\caption{A smooth matching neck.  The forcing is split in a fixed-width
middle interval.  The two decaying solutions are cut off over long
intervals, while their common tangent value is retained once.}
\label{local:neck-figure}
\end{figure}

The preceding proof also handles the controlled inhomogeneous collars
used later for ordinary continuation.  More precisely, suppose that,
after Hamiltonian transport of a limiting chord to an intersection,
the equation has an additional term $R_T(s,\tau,u)$ with bounded finite
strip derivatives, of size $O(r_l+r_r)$, and with differences from the
two partially broken equations $O(r_r)$ and $O(r_l)$ on the respective
halves.  Assume the differentiated versions in
Lemma~\ref{local:neck-coefficients} as well.  Each limiting end then
has a decaying forcing $O(e^{-\pi s})$.  The spectral inverse gives
decay at any smaller rate than the indicial gap and $\pi$; pregluing
has an exponentially small residual in a smaller weight, and the
same parametrix proof applies.  No nondecaying additive forcing is
permitted in this assertion.

\begin{proposition}[Local analytic chart contract]\label{local:charts}
For any prescribed finite integer $r$, choose sufficiently high Sobolev
order and sufficiently small positive weights.  Around each stable
map in the finite system of Proposition~\ref{local:finite}, the
component augmentations of Lemma~\ref{local:augmentation} give the
following local construction, also for $p_0$.

Allow finite-dimensional coefficient variables whose insertions are
supported on the normalized component cores, depend smoothly on
finitely many core jets and incidence markings, ignore the auxiliary
stabilizing marks, and have the core-transport and exponential
length-derivative properties of Lemma~\ref{local:neck-coefficients}.
At zero coefficients assume the fixed-domain augmentations above
are surjective.  Then the small pairs solving
\begin{equation}\label{local:thickened-equation}
 \bar\partial_Ju=I(u,\mathbf e)
\end{equation}
form local $C^r$ manifolds with corners modulo the finite map
stabilizers.  Corner coordinates are $q_e=1/T_e\geq0$, with $q_e=0$
meaning that node $e$ is retained.  Their normalized faces are exactly
the corresponding component solution spaces, matched over $L$ at
smooth nodes and over the fixed intersection at jumps.  Pure boundary
outputs are submersive on every stratum, also with the domain fixed.
Interior outputs need only be smooth.

These charts contain all sufficiently close solutions of
\eqref{local:thickened-equation}, including all partially broken ones.
Core observables, evaluation maps, and retrieval of transverse
stabilizing slices are $C^r$.  Changes between stabilizations retaining
the components and nodes at the comparison point are $C^r$ in these
coordinates.  Finally, outside sufficiently large fixed compact
interior cores, derivatives on hyperbolic unit vectors can be made
uniformly small on all these small-coefficient solutions.
\end{proposition}
\begin{proof}
We spell out both the nonlinear estimate and the length-derivative
estimate needed here.  Stabilizing slice equations are imposed with
marking positions free before division by disk automorphisms.  Their
transversality follows from the immersive incidence points, since
neither the equation nor the insertions uses those auxiliary marks.
Choose a finite collection $b$ of core kernel observables, including
coefficient and domain coordinates as necessary, which makes the
nodal linearization an isomorphism.  Near the nodal map they
parametrize its augmented solution stratum.  Preglue this nearby
nodal family with the prescribed $b$, cutting the decaying map parts
and retaining the matching values on $L$.

The residual of \eqref{local:thickened-equation}, including the
prescribed core observables, has norm at most
$C\sum_e e^{-cT_e}$ in a weight smaller than the component decay
rate.  On each core this follows from the smooth insertion transport
and Lemma~\ref{local:neck-coefficients}; in each middle neck it follows
from cutting exponentially decaying tails and the two one-sided
coefficient comparisons.  The same statement includes the controlled
additive errors described above, when present.

The nonlinear remainder satisfies, on a ball of fixed small radius,
\begin{equation}\label{local:quadratic}
 \|\mathcal N_T(\xi)-\mathcal N_T(\eta)\|_{m-1,\delta,T}
 \leq C\bigl(\|\xi\|_{m,\delta,T}+\|\eta\|_{m,\delta,T}\bigr)
                \|\xi-\eta\|_{m,\delta,T},
\end{equation}
with $C$ independent of the lengths.  Indeed ordinary Sobolev
multiplication on fixed-width cylinder charts has uniform constants;
$w_e\geq1$ puts the weight on the differentiated factor, while
Sobolev embedding bounds the undifferentiated one.  At a smooth
node use actual motions of the common point in $L$, not an affine
variation that violates the boundary condition.  In sheet-straightening
coordinates a constant such motion still solves the constant-map
equation.  Thus every nonlinear term on the neck involving the
constant has a derivative of a nonconstant part as a factor.  Its
weighted estimate is exactly the preceding product estimate, with
$|a_e|$ handled as a finite-dimensional variable.  The insertions
are on fixed cores, where ordinary high-order estimates apply.
Starting at a larger $m$ accommodates their finite jet dependence.

Apply the contraction theorem with the inverse in
Lemma~\ref{local:inverse} and the prescribed observables.  It gives a
unique small correction, with radius and estimates uniform in all
large lengths.  It works unchanged with an arbitrary subset of nodes
still broken, retaining their matching equations.  Conversely, a
nearby solution has the indicated stabilized domain by marked-domain
compactness.  Place its slice markings uniquely by incidence on the
retained cores.  Core closeness and Lemma~\ref{local:decay} put it in
the weighted implicit-function neighborhood: outside successively
larger cores the decaying tails are small in any strictly smaller
weight.  Hence the implicit construction captures every nearby
solution and has no extra local branch.  Lemma~\ref{local:augmentation}
and the version of the inverse with output prescribed prove
stratumwise output submersivity.  Kernel coordinates for this
statement may be chosen at fixed domain as in that lemma.

It remains to verify differentiability at simultaneous corners.
Identify domains under a change of $T_e$ by fixing vertex coordinates
on both sides and stretching only a middle part of neck $e$.
Trivialize compact component deformations while keeping their disjoint
nodal coordinate germs fixed.  Fix a finite range
\[
 m_0>m_1>\cdots>m_r>3,\qquad
 \delta_0>\delta_1>\cdots>\delta_r>0,
\]
with sufficiently large initial order.  The nodal and glued inverses
are bounded in this entire range.

We record the weight estimate used for a middle stretch.
Let $\chi_e$ be supported where
$\min(s,T_e-s)\geq\theta T_e$, for a fixed $0<\theta<1/2$, with
uniformly bounded strip derivatives up to the order in use.
For two successive weights $\delta'>\delta''$, using the weighted
target norm on the left,
\[
 \|\chi_e D_{\rm strip}\xi\|_{m-1,\delta'',T}
 \leq C e^{-\theta(\delta'-\delta'')T_e}
                   \|\xi\|_{m,\delta',T}.
\]
Indeed, the ratio of the two weights on this support is bounded by a
uniform multiple of
$e^{-\theta(\delta'-\delta'')T_e}$.  Here $D_{\rm strip}$ differentiates
components in the sheet-straightening coordinates: at a smooth match,
$D_{\rm strip}(a_e+\xi^{\rm dec})=D_{\rm strip}\xi^{\rm dec}$.
The estimate therefore applies to lower parameter variations as well
as to the map remainder.  A polynomial in $T_e$ multiplying the
left side is absorbed by reducing the exponential rate.
Differentiating the domain identification produces these middle
strip-derivative terms.  Its remaining zero-order terms involving a
movable constant contain a strip derivative of the map or a lower
variation, and obey the same weighted product estimate.
For a stretch on another neck $f$, the loss of weight there absorbs
its polynomial in $T_f$, uniformly in all the other lengths.

Differentiating the equation with
$b$ fixed uses the same inverse.  Explicit length derivatives of its
core coefficients and insertions are exponentially small.  On neck
$e$, stretching derivatives multiply a positive strip derivative of
the map, exponentially small near the middle in the next smaller
weight.  On any other neck, the differentiated harmonic coefficient
has the factor $e^{-\pi T_e}$ established in
Lemma~\ref{local:neck-coefficients}; it multiplies a strip derivative
of the map.  Constants occur only at smooth matches, where the
coefficient is $J_*$ and every relevant zero-order term has a factor
$du$.  These observations give, in the next norm of the finite range,
\[
 \|\partial_{T_e}u_T\|_{m_1,\delta_1,T}
           \leq C e^{-cT_e}.
\]
Here derivatives mean the covariant derivative under the just
specified core and middle-stretch identifications, including the
movable constants.

The same argument iterates.  At order $j$ the unknown highest
variation is acted on by the same augmented linearization; its
right-hand side is a sum of products of lower variations, differentiated
coefficients, and positive strip derivatives.  The uniform product
estimate, with a finite loss of Sobolev order, bounds every such term.
A middle-stretch derivative may contribute a polynomial in its length;
the difference between successive weights absorbs this polynomial
against the exponential tail.  If a mixed derivative contains
$\partial_{T_e}$, apply this induction with that derivative first.
Every subsequent tangential or length differentiation preserves an
exponential factor in $T_e$, in a weaker norm.  Thus for every
multi-index $\alpha$ with $1\leq|\alpha|\leq r$ and every $e$ in its
support, and for the allowed remaining parameter derivatives $D_b^\gamma$,
\begin{equation}\label{local:length-derivatives}
 \|D_b^\gamma\partial_T^\alpha u_T\|_{m_r,\delta_r,T}
              \leq C_{\alpha,\gamma}e^{-c_rT_e},
 \qquad |\alpha|+|\gamma|\leq r.
\end{equation}
Constants are uniform in the other large lengths and in which of
the other nodes are held broken.  This also bounds the restrictions
to every retained core and their movable limiting values.  Using
an index $e$ for which $T_e$ is largest among the differentiated
lengths absorbs every polynomial introduced by
$\partial_{q_e}=-T_e^2\partial_{T_e}$.

The limit as $T_e\to\infty$ is the corresponding partially broken
solution, by matching and implicit uniqueness.  Integrating the
exponential derivative bounds to that face, also after the allowed
tangential differentiations, proves $C^r$ extension in all $q_e$.
At a smooth node held broken, the estimate initially controls the
value and parameter variations uniformly on the punctured end.
In a smaller unpunctured neighborhood the equation has fixed $J_*$
and no obstruction support.  Removable-boundary-point regularity and
elliptic estimates, also for the differentiated equations, upgrade
these bounds to all needed evaluations at the node.  This completes
the differentiability proof for evaluations and core observables.

Two systems of auxiliary slices are recovered smoothly on retained
cores by their transverse incidence equations.  Compare them, if
necessary, through a third disjoint stabilization.  At the comparison
point all components are retained in each description; no chain of
nodes is contracted.  The real plumbing parameters for the same
node consequently satisfy
\[
 \rho'_e=v_e\rho_e,\qquad v_e>0,
\]
where $v_e$ is a smooth unit in the marked-domain coordinates.  After
retrieving the slices it is $C^r$ in the chart variables.  Therefore
\[
 T'_e=T_e-\pi^{-1}\log v_e,\qquad
 q'_e=\frac{q_e}{1-\pi^{-1}q_e\log v_e},
\]
which is $C^r$ also at $q_e=0$.  This argument concerns retained real
boundary nodes only.  The angular interior-node problem addressed
in the corrigendum \cite[Section~1]{fooo-corr} never arises here.
No compatibility with forgetting actual input markings is asserted
or needed.

Finally all insertions stay on compact interior cores.  Near ordinary
boundary arcs, the hyperbolic metric makes a bounded smooth derivative
small on its unit vectors as the boundary is approached.  In jump
ends and necks, hyperbolic unit vectors are bounded in strip
coordinates and \eqref{local:two-sided-decay} applies.  Choose the
cores large enough and then shrink the coefficient and implicit
neighborhoods.  These observations give the stated uniform smallness
off the cores, including for the nonzero-coefficient solutions used
in the next section.  Finite stabilizers act on this entire construction;
passing to their quotients gives the announced local charts.
\end{proof}
